% Fichier généré automatiquement par tools/generate_corrections.py.
% Source : SLCI/SLCI-03-SchemaBlocs/48_Quille/48_Quille.tex
% Statut : complet (3/3 question(s) renseignée(s), 0 reprise(s)).
\begin{corrigebox}[Corrigé extrait de la banque]
\CorrigeQuestion{1}
D'une part, on transforme les équations dans le domaine de Laplace : 
$Q(p)=S p X(p)+\dfrac{V}{2B} p \Sigma(p)$ et
$Mp^2 X(p) = S \Sigma(p) - kX(p)-\lambda p X(p) - F_R(p)$.

En utilisant le schéma-blocs, on a $\Sigma(p)=A_2\left(A_1Q(p)-X(p)\right) = A_1A_2Q(p)-A_2X(p)$.

Par ailleurs $\Sigma(p)=\dfrac{Q(p)-S p X(p)}{\dfrac{V}{2B} p}= Q(p)\dfrac{2B}{Vp}-  X(p)  \dfrac{S2B}{V} $. On a donc $A_2 = \dfrac{S2B}{V} $, $A_1 A_2 = \dfrac{2B}{Vp}$ soit $A_1  = \dfrac{2B}{Vp}\dfrac{V}{S2B}= \dfrac{1}{Sp}$. 

On a aussi $X(p)=A_4\left(-F_R(p)+A_3\Sigma(p)\right) =-A_4F_R(p)+A_3A_4\Sigma(p)$. Par ailleurs,
$X(p) \left(Mp^2  +\lambda p  + k\right)= S \Sigma(p) - F_R(p) \Leftrightarrow X(p) =  \dfrac{S \Sigma(p)}{Mp^2  +\lambda p  + k}-\dfrac{F_R(p)}{Mp^2  +\lambda p  + k}$. On a donc : $A_4 = \dfrac{1}{Mp^2  +\lambda p  + k}$ et $A_3 = S$.

Au final,  $A_1=\dfrac{1}{Sp}$,  $A_2 = \dfrac{S2B}{V} $,  $A_3 = S$  et $A_4 = \dfrac{1}{Mp^2  +\lambda p  + k}$.
\CorrigeQuestion{2}
\textbf{Méthode 1 : Utilisation des relations précédentes}
On a $X(p)=\left(H_1Q(p)-F_R(p)\right)H_2(p)$. 

Par ailleurs, on a vu que $X(p)=A_4\left(-F_R(p)+A_3\Sigma(p)\right) $ et $\Sigma(p)=A_2\left(A_1Q(p)-X(p)\right)$. 

On a donc $X(p)=A_4\left(-F_R(p)+A_3  A_2\left(A_1Q(p)-X(p)\right)\right) $ $ \Leftrightarrow X(p)\left(1+A_2A_3A_4 \right)=A_4\left(-F_R(p)+A_3  A_2A_1Q(p)\right) $. On a donc 
$H_1(p)=A_1  A_2A_3$ et $H_2 = \dfrac{A_4}{1+ A_2A_3A_4 }$.

\textbf{Méthode 2 : Lecture directe du schéma-blocs}
Revient à utiliser la méthode précédente. 

\textbf{Méthode 3 : Algèbre de schéma-blocs}
Le schéma-blocs proposé est équivalent au schéma suivant. 

\footnotesize
\par\medskip\begin{center}\captionsetup{type=figure}
\begin{tikzpicture}
\sbEntree{E}

\sbBloc[3]{b0}{$A_1 A_2 A_3$}{E}
    \sbRelier[$Q(p)$]{E}{b0}

\sbComph{c2}{b0}
    \sbRelier{b0}{c2}
    
\sbComp{c1}{c2}
    \sbRelier{c2}{c1}

\sbBloc{b1}{$A_4$}{c1}
    \sbRelier{c1}{b1}
    

\sbSortie[4]{S}{b1}
    \sbRelier{b1}{S}
    \sbNomLien[0.8]{S}{$X(p)$}
      
      
\sbDecaleNoeudy[4]{S}{U}
\sbDecaleNoeudx[-2]{U}{U2}
\sbBlocr{r1}{$A_2 A_3$}{U2}

\sbRelieryx{b1-S}{r1}
\sbRelierxy{r1}{c1}

\draw [latex-] (c2) --++ (0,1) node[left] {$F_R(p)$};

\end{tikzpicture}
\end{center}\par\medskip
\normalsize

On retrouve le même résultat que précédemment. 

$A_1=\dfrac{1}{Sp}$,  $A_2 = \dfrac{S2B}{V} $,  $A_3 = S$  et $A_4 = \dfrac{1}{Mp^2  +\lambda p  + k}$.

En faisant le calcul on obtient : 
$H_1(p)=\dfrac{2BS}{pV}  $ et $H_2 = \dfrac{\dfrac{1}{Mp^2  +\lambda p  + k}}{1+ \dfrac{2BS^2}{V}\dfrac{1}{Mp^2  +\lambda p  + k} }$  $= \dfrac{1}{Mp^2  +\lambda p  + k+ \dfrac{2BS^2}{V} }$.
\CorrigeQuestion{3}
Dans ce cas, $\dfrac{X(p)}{Q(p)}=H_1(p)H_2(p)=\dfrac{2BS}{p\left(MVp^2  +\lambda pV  + kV+ 2BS^2\right) }$.

\begin{enumerate}
\item $A_1=\dfrac{1}{Sp}$,  $A_2 = \dfrac{S2B}{V} $,  $A_3 = S$  et $A_4 = \dfrac{1}{Mp^2  +\lambda p  + k}$.
\item $H_1(p)=A_1  A_2A_3$ et $H_2 = \dfrac{A_4}{1+ A_2A_3A_4 }$.
\item $\dfrac{X(p)}{Q(p)}=\dfrac{2BS}{p\left(MVp^2  +\lambda pV  + kV+ 2BS^2\right) }$.
\end{enumerate}
\end{corrigebox}
